% =========================================================================
% CAPÍTULO III: METODOLOGÍA DE LA INVESTIGACIÓN
% =========================================================================
\clearpage
\chapter{Metodología de la investigación}
\label{sec:metodologia}

\section{Enfoque, tipo y diseño de investigación}
\subsection{Enfoque de investigación}
El estudio tendrá un enfoque cuantitativo: comparará configuraciones de software mediante indicadores numéricos de recuperación, calidad de respuesta y eficiencia computacional. La unidad de evaluación será la respuesta producida para una consulta bajo una configuración determinada. No se medirán efectos educativos en estudiantes ni impactos sociales.

\subsection{Tipo de investigación}
La investigación será aplicada y tecnológica porque diseñará, implementará y evaluará un artefacto de software para una necesidad delimitada. El proceso seguirá las actividades de diseño, construcción y evaluación de artefactos de la investigación en ciencias del diseño (DSR, de \textit{Design Science Research}) \parencite{hevner2004design}.

\subsection{Diseño de investigación}
Se realizará un experimento computacional comparativo con medidas repetidas por consulta. Las condiciones de generación serán: (a) SLM base; (b) el mismo SLM con recuperación documental local (SLM--RAG); y (c) SLM--RAG con adaptadores QLoRA, únicamente si el entrenamiento y la inferencia adaptada son viables. Para la hipótesis de recuperación se contrastará, por separado, el índice semántico con una búsqueda léxica de referencia. En cada comparación se mantendrán constantes las consultas, los documentos disponibles, el hardware y los parámetros pertinentes. La comparación directa para estimar el aporte de QLoRA será SLM--RAG frente a SLM--RAG--QLoRA; la comparación con SLM base describirá el efecto conjunto de recuperación y adaptación.

\section{Variables y operacionalización}
\subsection{Identificación de variables}
\begin{itemize}
    \item \textbf{Variable independiente:} configuración del sistema: SLM base, SLM--RAG o SLM--RAG--QLoRA, cuando esta última pueda ejecutarse.
    \item \textbf{Variables dependientes:} recuperación de fragmentos pertinentes, similitud de las respuestas con referencias validadas, fidelidad al contexto y eficiencia computacional.
    \item \textbf{Condiciones de control:} versión del modelo y del tokenizador, cuantización, colección documental, consultas, índice, plantilla de instrucciones, parámetros de generación, equipo y protocolo de medición. Los valores concretos se registrarán antes de la comparación.
\end{itemize}

\subsection{Matriz de operacionalización de variables}
La Tabla~\ref{tab:operacionalizacion_variables} describe los indicadores y el procedimiento previsto para obtenerlos. Las métricas automáticas de similitud se interpretarán como aproximaciones respecto de una referencia, no como una medida completa de corrección o fluidez.

\begin{table}[H]
\centering
\begin{threeparttable}
\caption{Matriz de operacionalización de variables}
\label{tab:operacionalizacion_variables}
\small
\begin{tabularx}{\textwidth}{@{}p{2.25cm}p{3.25cm}p{3.65cm}X@{}}
\toprule
\textbf{Variable} & \textbf{Operacionalización} & \textbf{Indicadores} & \textbf{Escala e instrumento} \\
\midrule
Configuración (independiente) & Condición aplicada a cada consulta: SLM, SLM--RAG o SLM--RAG--QLoRA, si es viable. & Identificador de configuración; modelo, índice y fragmentos usados. & Categórica; configuración versionada y ficha del Anexo 2. \\
\addlinespace
Recuperación & Pertinencia de los primeros $k$ fragmentos para cada consulta. & Precisión en los primeros $k$ resultados (P@k), índice semántico frente a búsqueda léxica. & Numérica; banco de consultas, etiquetas de relevancia y registro del recuperador. \\
\addlinespace
Respuesta y fidelidad & Comparación de cada salida con una respuesta de referencia y con la evidencia entregada al modelo. & chrF++; BLEU; fidelidad y recuperación de contexto, si se valida el evaluador automático. & Numérica; SacreBLEU, implementación documentada de las métricas RAG y revisión con la rúbrica del Anexo 3. \\
\addlinespace
Eficiencia & Costo de ejecución por consulta y configuración en el mismo equipo. & Latencia; tokens generados por segundo; máximos de RAM y VRAM; errores de ejecución. & Numérica; temporizador, telemetría del sistema y ficha del Anexo 2. \\
\bottomrule
\end{tabularx}
\begin{tablenotes}
\small
\item \textit{Nota.} Se registrarán las versiones de herramientas, los parámetros, las reglas de normalización y los valores de $k$. Las métricas RAG que dependan de un modelo evaluador solo se usarán de forma confirmatoria si se verifica su capacidad para evaluar respuestas en Quechua Collao.
\end{tablenotes}
\end{threeparttable}
\end{table}

\section{Unidades de análisis y selección de consultas}
\subsection{Unidad de análisis}
Para la evaluación de respuestas, la unidad de análisis será el par formado por una consulta y la respuesta generada bajo una configuración, junto con los fragmentos recuperados y el registro de ejecución. Para la comparación de recuperadores, la unidad será la consulta con sus etiquetas de pertinencia. Cada consulta se mantendrá emparejada entre las condiciones que se comparen.

\subsection{Población y muestra}
La población conceptual comprende consultas textuales que puedan responderse con fuentes autorizadas identificadas como Quechua Collao o español. La factibilidad propone un banco intencional de 100 consultas patrón; se conserva como meta de planificación, sujeta al inventario de fuentes y a un piloto. No proviene aún de un cálculo de potencia estadística. Se incluirán consultas con respuesta sustentable en una fuente rastreable y se excluirán las que requieran información ausente o ambigua en el corpus. Si se ejecuta QLoRA, las consultas de evaluación y sus respuestas de referencia se reservarán fuera de los ejemplos de ajuste; el tratamiento de documentos indexados y coincidencias con ejemplos de entrenamiento seguirá el procedimiento de la subsección siguiente.

\subsection{Preparación del corpus y prevención de contaminación}
Se inventariarán las fuentes propuestas en el estudio de factibilidad —corpus bilingües, recursos léxicos y textos abiertos— y se registrarán su procedencia, licencia, idioma y variedad. Los materiales obtenidos mediante reconocimiento óptico de caracteres (OCR) se revisarán antes de indexarlos; se conservarán los datos de procedencia y las correcciones aplicadas. Los documentos se limpiarán y segmentarán. La fuente del proyecto propone fragmentos de 250--400 tokens y recuperar tres fragmentos por consulta; ambos valores se tratarán como puntos de partida para el piloto, no como parámetros ya validados, y se fijarán antes de la evaluación comparativa.

La propuesta fuente también contempla pares de pregunta y respuesta generados con asistencia de modelos de lenguaje. Si se emplean, se identificarán como datos sintéticos y no se considerarán referencias correctas por el hecho de haber sido generados. Se contrastarán con fuentes trazables y se someterán a revisión lingüística competente antes de servir como ejemplos de ajuste o referencias de evaluación. Las consultas de evaluación y sus respuestas de referencia se mantendrán fuera del ajuste QLoRA. En cambio, los documentos con la evidencia necesaria para responderlas permanecerán disponibles en la base de conocimiento de las condiciones RAG, con el mismo contenido para SLM--RAG y SLM--RAG--QLoRA; se registrarán coincidencias textuales relevantes entre datos de ajuste y documentos indexados. Así se evita excluir del índice la evidencia que el sistema debe recuperar.

\section{Técnicas e instrumentos de recolección}
La técnica principal será la ejecución controlada del prototipo con captura automatizada de consultas, respuestas, fragmentos recuperados, latencia, consumo de memoria y errores. Los instrumentos serán el inventario de fuentes, el banco de consultas del Anexo 4, la ficha de ejecución del Anexo 2, los scripts versionados de métricas y la rúbrica de revisión lingüística y factual del Anexo 3.

Antes de calcular métricas con referencias, se documentará para cada consulta su idioma y variedad, fuente, respuesta de referencia y evidencia pertinente. Una persona competente en Quechua Collao deberá revisar las referencias destinadas a medir similitud lingüística. Si esa revisión no se consigue, chrF++ y BLEU se informarán como similitud automática con referencias no validadas, y no como prueba suficiente de calidad lingüística. Una métrica automática de fidelidad que requiera un modelo juez se considerará exploratoria hasta comprobar, en un piloto documentado, que el juez puede evaluar el idioma y la variedad del estudio.

\section{Procedimiento metodológico}
El procedimiento conserva las tres fases y la correspondencia con los objetivos específicos de la presentación de objetivos:
\begin{enumerate}
    \item \textbf{Fase 1 — Datos y recuperación (OE1):} inventariar las fuentes y sus licencias; etiquetar idioma, variedad y procedencia; revisar y depurar los documentos; pilotear la segmentación y el número de fragmentos recuperados; crear el índice semántico y una búsqueda léxica de referencia; y preparar consultas con sus fuentes y fragmentos pertinentes identificados.
    \item \textbf{Fase 2 — Núcleo técnico y adaptación (OE2):} comprobar la inferencia local de Qwen2.5-3B-Instruct y del recuperador; fijar y registrar las versiones, el formato de cuantización, las instrucciones y los parámetros de generación; integrar SLM y RAG; y realizar una prueba técnica preliminar de ajuste, exportación, carga e inferencia con QLoRA para determinar si puede ejecutarse la condición E3. Esta comprobación forma parte del procedimiento metodológico; el análisis global de factibilidad técnica, económica y operativa, así como de viabilidad y sostenibilidad, se desarrollará por separado en el capítulo IV. La fuente propone Q4\_K\_M/GGUF con \texttt{llama.cpp} para inferencia y contempla \texttt{bitsandbytes}/PEFT para el ajuste; son etapas y formatos distintos, y se verificará que la exportación y carga del adaptador sean compatibles. En este equipo se prevé evaluar WSL2 con Ubuntu para \texttt{bitsandbytes}. Si el ajuste o la inferencia adaptada fallan o no son reproducibles, se registrará el motivo y no se ejecutará la condición integrada.
    \item \textbf{Fase 3 — Evaluación y análisis (OE3):} ejecutar el mismo banco bajo cada condición viable; registrar entradas, salidas, fragmentos y telemetría; calcular las métricas definidas; inspeccionar errores con criterios comunes; y conservar los scripts, parámetros y registros suficientes para repetir el procedimiento.
\end{enumerate}

\section{Protocolo de evaluación}
Antes de la corrida comparativa se cerrará un protocolo que identifique el modelo y su revisión, cuantización de inferencia y de ajuste, tokenizador, motor de inferencia, modelo de embeddings, segmentación documental, valor de $k$, plantilla de instrucciones y parámetros de decodificación. Se ejecutará la evaluación sin conexión y con los mismos recursos de hardware. Se registrarán el sistema operativo y las versiones de software. Las consultas de evaluación no se usarán para seleccionar parámetros ni ajustar QLoRA. El número de repeticiones para mediciones de rendimiento se fijará tras el piloto y se aplicará igual a todas las condiciones; las respuestas utilizadas para métricas de texto se generarán con parámetros de decodificación constantes.

\subsection{Métricas de calidad lingüística}
Se calcularán chrF++ y BLEU mediante SacreBLEU, comparando las salidas con referencias documentadas y, cuando sea posible, revisadas por una persona competente en la variedad. Se conservarán la versión de SacreBLEU, su firma de configuración, la tokenización, la normalización, el idioma de salida y la cantidad de consultas \parencite{popovic-2017-chrf,Post2018}. Estas métricas cuantifican similitud con referencias y no se interpretarán por sí solas como comprensión, corrección semántica o fluidez.

\subsection{Métricas de recuperación y fidelidad}
Se calculará la precisión en los primeros $k$ resultados (P@k) usando etiquetas de pertinencia definidas antes de evaluar y el mismo $k$ para la búsqueda léxica y semántica. Para las respuestas se registrará si sus afirmaciones están respaldadas por los fragmentos proporcionados. Las métricas de RAGAS seleccionadas podrán complementar esa revisión solo si la versión y el modelo evaluador se identifican y el piloto confirma su uso pertinente con Quechua Collao; de lo contrario, sus puntuaciones no se presentarán como evidencia confirmatoria \parencite{Ragas2024}.

\subsection{Métricas de eficiencia}
Por consulta y condición se medirán latencia, tokens generados por segundo y máximos de memoria principal (RAM) y memoria gráfica (VRAM). Se informarán medidas descriptivas, dispersión, fallos y datos faltantes. Cuando la instrumentación lo permita, se separarán los tiempos de recuperación, preparación del contexto e inferencia. Las mediciones de rendimiento seguirán el mismo procedimiento de calentamiento y repetición en cada condición, definido y registrado antes de la comparación.

\section{Métodos de análisis de datos}
Se calcularán estadísticos descriptivos por configuración y métrica. Para HE1 se comparará P@k del recuperador semántico con la búsqueda léxica sobre las mismas consultas. Para HE2, si QLoRA es viable, la comparación principal será SLM--RAG frente a SLM--RAG--QLoRA, manteniendo los demás parámetros. Para HE3 se describirán y, si los datos emparejados y los supuestos lo permiten, se contrastarán las tres configuraciones. Friedman se usará para tres condiciones relacionadas y Wilcoxon de rangos con signo para comparaciones pareadas; las comparaciones múltiples se ajustarán mediante Holm. Se comprobará que la escala y los datos disponibles permitan aplicar estas pruebas; si no, el análisis será descriptivo y se informará la limitación. Se fijará $\alpha=0{,}05$ antes del análisis y se informarán tamaños de efecto, intervalos de confianza cuando sean estimables, datos faltantes y ajustes por comparaciones múltiples. Si QLoRA no es viable, HE2 y las comparaciones de HE3 que la involucren se declararán no evaluadas; no se inferirá superioridad de una condición que no se ejecutó.

\section{Aspectos éticos y de gestión de datos}
Se incorporarán documentos con procedencia y condiciones de uso verificables, se citarán las fuentes y no se redistribuirán contenidos protegidos sin autorización. Las consultas y bitácoras no incluirán datos personales. Los datos generados artificialmente se identificarán y separarán de los textos fuente. Se conservarán las etiquetas de Quechua Collao y no se mezclarán variedades sin identificación explícita. Si la validación lingüística requiere participación de personas, se obtendrán las autorizaciones y el consentimiento informado que correspondan antes de recoger sus evaluaciones.
